\begin{abstract}
Vision--language--action (VLA) systems map language and multimodal observations
to continuous robot actions.  At deployment, an online action chunk crosses
from model output into the execution stack without inherent evidence that it
is admissible for the current task, that later software preserves it, or that
its realization is qualified for the current robot state.  Existing defenses
typically stop at model input or candidate selection, or begin from an already
trusted command, program, or reference trajectory.

We present \system, an inline consumer-side reference monitor that makes the
exact source \actionblock the protected object.  Under an explicit trusted-task
and observation assumption, \lone attaches a checker-relative assessment to
the block.  \ltwoa carries its identity through a fresh one-use authorization,
ordered dispatch, receipts, effects, and task-phase advance.  When a registered
joint boundary triggers, \ltwob qualifies at most two temporary virtual-guard
configurations without changing or resampling the source action.  We formalize
the composition as $\mathsf{ProofAligned}$---checker-relative eligibility,
transaction alignment, and conditional containment---and use Lean to
machine-check its finite authorization, ordered-action, receipt, evidence, and
phase-transition relations.

Using SABER on a complete 120-unit attack-evaluation protocol, we measure 39 of
86 clean-eligible units as new risk transitions (45.35\%, 95\% base-pair
cluster-bootstrap CI: [32.93\%, 57.78\%]).  A separate paired four-arm study
contains 18 task/initialization pairs and 144 episodes.  Under attack, VLA-only
and L2-only succeed on 11/18 tasks, whereas L1-only and Dual succeed on 13/18;
their joint-limit violation-episode counts are 4/18, 1/18, 0/18, and 0/18.
Screening reaches 39.79\,ms maximum and 18.30\,ms p95, with no 100\,ms miss.
The empirical studies run in LIBERO-Safety simulation; they support the stated
checker-relative and transaction properties for the evaluated configuration,
not complete semantic correctness, hardware attestation, hard-real-time
behavior, or general robot safety.
\end{abstract}
